% FILE: 04_implementation_surface.tex
% Project: otel-weaver
% Role: telemetry-to-process-evidence-ingestion-layer

\section{Implementation Surface}
\label{sec:otel-weaver-impl}

\subsection{Source Files}
\label{sec:otel-weaver-impl-sources}

The implementation surface of \texttt{otel-weaver} spans five experiment subdirectories, six
doctrine files, five YAML mapping tables, and supporting \texttt{ggen} and \texttt{intel}
artifacts. The following experiment-level source files are the primary loci of runnable
specification:

\paragraph{exp-001: Custom PI Weaver Registry.}
Defined in \texttt{experiments/exp-001-custom-pi-weaver-registry/README.md}. This experiment
specifies the \texttt{process.pi} semantic convention group as a custom OTel Weaver registry.
The registry defines six required attribute fields: \texttt{instance\_id} (process instance
identifier), \texttt{activity.name} (declared activity label), \texttt{lifecycle} (lifecycle
state string), token state fields \texttt{before} and \texttt{after}, \texttt{witness.id}
(witness identifier), and \texttt{witness.hash} (BLAKE3 64-character digest). The registry
is the feedstock schema; passing it is necessary but not sufficient for process-evidence status.

\paragraph{exp-002: Weaver Diff to PI Residual.}
Defined in \texttt{experiments/exp-002-weaver-diff-to-pi-residual/README.md} with an accompanying
\texttt{Cargo.toml} declaring dependencies on \texttt{serde} and \texttt{blake3}. The Rust
implementation in \texttt{src/bridge.rs} specifies the \texttt{BridgeRx} struct and its
core property: $R_\Delta = |f_{S_1}(L_2, M) - f_{S_2}(L_2, M)| = 0$ after bridge translation.
The PI conformance residual formula $R_\Delta = |f_{S_1}(L_2, M) - f_{S_2}(L_2, M)|$ is the
formal measurement of error introduced by schema mismatch between two Weaver schema versions
$S_1$ and $S_2$ when replaying event log $L_2$ against process model $M$.

\paragraph{exp-003: Live-Check to Refusal.}
Defined in \texttt{experiments/exp-003-live-check-to-refusal/README.md} with a \texttt{Cargo.toml}
declaring a path dependency on \texttt{wasm4pm-compat} at \texttt{../../../../wasm4pm-compat}.
The Rust implementation in \texttt{src/validator.rs} specifies the \texttt{Admission<T,\,W>}
gate function and the \texttt{Refusal} struct. Two inline test cases are fully specified:
\texttt{test\_successful\_admission} (verifies that a well-formed \texttt{process.pi} span
produces an \texttt{Admission} wrapping an \texttt{ActivityWitness}) and
\texttt{test\_refusal\_missing\_id} (verifies that a span missing \texttt{instance\_id}
produces a \texttt{Refusal} with \texttt{RefusalCode::MissingRequiredField} and a BLAKE3
digest of the rejected payload).

\paragraph{exp-004: Registry to wasm4pm-compat Witness.}
Defined in \texttt{experiments/exp-004-registry-to-wasm4pm-compat-witness/README.md}. This
experiment specifies the \texttt{ActivityWitness} lattice struct and its code-materialization
path from a resolved OTel Weaver registry JSON. The \texttt{Lattice} trait implementation
uses witness hash emptiness as the partial order criterion. The codegen path produces
\texttt{ActivityWitness} Rust source as output of the Weaver \texttt{generate} command applied
to the \texttt{process.pi} registry.

\paragraph{exp-005: Collector to PI Intake.}
Defined in \texttt{experiments/exp-005-collector-to-pi-intake/README.md}. This experiment
specifies the OTel Collector YAML configuration that routes spans through the ingestion pipeline.
The configuration declares receiver, processor, and exporter stages that enforce the Ingestion
Covenant: Admit-Don't-Assume (no span assumed to be evidence without gate passage), Loss Reports
(every projection declares its \texttt{OtelToOcelLossPolicy}), and Witness Seal (every admitted
payload carries an \texttt{ActivityWitness}).

\subsection{Commands}
\label{sec:otel-weaver-impl-commands}

The \texttt{intel/weaver-command-map.yaml} file maps the full OTel Weaver command surface to
ingestion-layer roles. The commands relevant to this project's implementation are:

\begin{description}
  \item[\texttt{weaver registry check}] Validates a registry definition against its schema.
    Used in experiment \texttt{exp-001} to verify the \texttt{process.pi} registry is
    well-formed before it is referenced by the admission gate.
  \item[\texttt{weaver registry live-check}] Real-time feedstock validation gateway. The
    primary command in \texttt{exp-003}: incoming spans are validated against the
    \texttt{process.pi} registry in real time. A passing live-check is a precondition for
    the admission gate but does not constitute admission.
  \item[\texttt{weaver registry diff}] Computes $\Delta R = R_{\text{new}} \ominus R_{\text{old}}$
    between two registry versions. Used in \texttt{exp-002} as the input to \texttt{BridgeRx}.
  \item[\texttt{weaver registry generate}] Materializes Rust source from a resolved registry
    JSON. Used in \texttt{exp-004} to manufacture the \texttt{ActivityWitness} struct.
  \item[\texttt{weaver registry stats}] Reports registry coverage and attribute population
    statistics. Used in audit scripts to verify \texttt{process.pi} field coverage across
    ingested span samples.
  \item[\texttt{weaver registry infer}] Infers a registry schema from observed span samples.
    Used in the gap-analysis phase to identify \texttt{process.pi} fields absent from
    production telemetry streams.
\end{description}

Additional commands documented in \texttt{weaver-command-map.yaml} include \texttt{update-markdown},
\texttt{json-schema}, \texttt{package}, and \texttt{mcp}, which support registry documentation
and distribution but are not primary ingestion-layer mechanisms.

\subsection{Data and Ontology Surfaces}
\label{sec:otel-weaver-impl-data}

The project's data surface is expressed entirely through five YAML mapping files and the
\texttt{process.pi} registry definition. There is no RDF/SPARQL layer; this is an open question
noted in Section~\ref{sec:otel-weaver-open-rdf}.

\begin{description}
  \item[\texttt{otel-to-pi-witness-map.yaml}] Maps OTel span attribute paths to
    \texttt{ActivityWitness} struct fields. Status: ACTIVE.
  \item[\texttt{otel-signal-to-process-evidence-map.yaml}] Maps OTel signal types
    (span, metric, log) to process-evidence classification (feedstock, inadmissible,
    conditional). Status: ACTIVE.
  \item[\texttt{registry-diff-to-pi-drift-map.yaml}] Maps Weaver registry diff categories
    to the explicit verdict that each category is \emph{not} process drift. Status: ACTIVE.
  \item[\texttt{live-check-finding-to-refusal-map.yaml}] Maps Weaver live-check failure
    codes to \texttt{RefusalCode} enum variants with \texttt{violated\_rule} references.
    Status: ACTIVE.
  \item[\texttt{schema-url-to-receipt-context-map.yaml}] Maps OTel schema URLs to receipt
    context metadata required for BLAKE3 chain construction. Status: ACTIVE.
\end{description}

The \texttt{intel/weaver-registry-model.yaml} file provides a structural model of the Weaver
registry format, used as a reference during \texttt{ActivityWitness} codegen.

\subsection{Templates and Rendered Artifacts}
\label{sec:otel-weaver-impl-templates}

Three \texttt{ggen} template files are referenced in the project's manifest chain:
\texttt{pi-live-check-intake.rs.ggen}, \texttt{pi-witness-map.rs.ggen}, and
\texttt{pi-otel-constants.rs.ggen}. These templates specify the Rust source to be manufactured
by \texttt{ggen} actuation. At time of writing, no \texttt{generated/} directory exists in the
repository tree, indicating that the templates have been authored but \texttt{ggen} actuation
has not materialized the output files to disk. This constitutes a deferred gate (see
Section~\ref{sec:otel-weaver-deferred-gates}). The two \texttt{ggen} manifests
(\texttt{otel-weaver-source.manifest.yaml} and \texttt{pi-telemetry-bridge.manifest.yaml})
reference BLAKE3 receipt files at \texttt{receipts/otel-weaver-source-receipt.json} and
\texttt{receipts/pi-telemetry-bridge-receipt.json}; these receipt files were not verified on
disk at the time of evidence extraction.
